\section{Discussion, Limitations, and Future Work}\label{sec:disc}


\subsection{Synthesis and Fleet Engineering Selection Matrix}
Stepping back across RQ1--RQ4, resolving head-on single-lane encounters separates cleanly into two choices. The first is selecting a \emph{decision mechanism} based on fleet diversity, travel-time estimation accuracy, and whether robots belong to a single owner or competing vendors. The second is configuring the \emph{execution protocol} for wireless reliability, map completeness, and multi-robot geometry. Table~\ref{tab:guide} summarizes these quantitative findings into a practical selection guide for AMR system designers.

\begin{table*}[!t]\centering
\caption{Practitioner policy and protocol selection matrix distilled from $>50{,}000$ 2D and 496 ROS~2 / Gazebo episodes.}
\label{tab:guide}
{\footnotesize\renewcommand{\arraystretch}{1.22}\setlength{\tabcolsep}{6pt}
\begin{tabular}{@{}>{\raggedright\arraybackslash}p{0.195\textwidth}>{\raggedright\arraybackslash}p{0.205\textwidth}>{\raggedright\arraybackslash}p{0.240\textwidth}>{\raggedright\arraybackslash}p{0.290\textwidth}@{}}\toprule
Deployment Regime & Recommended Policy & Protocol Configuration & Empirical Basis \& Caveats \\\midrule
Homogeneous cooperative fleet; uncertain kinematics ($\sigma \ge 0.5$) & \texttt{rules\_proximity} (single-bay aisles) or \texttt{fixed\_priority}/\texttt{rules\_default} (multi-bay aisles); $\lambda=0.1$ credit where delay estimates exist & Periodic \texttt{STATE}/\texttt{BID}; \texttt{share\_map=true}; \texttt{commit} optional & Single bay: halves ID-priority excess ($3.1$ vs.\ $6.4$\,s) and matches noisy \cb\ at $\sigma\ge 0.5$; two bays: $5.6$ vs.\ $1.1$\,s for fixed priority. Credit rule (held-out $\lambda$) reaches $J=0.992$ for $+3.8\%$ delay ($+11.5\%$ at $\sigma=0.5$) \\\addlinespace[4.5pt]
Heterogeneous cooperative fleet (urgent / loaded / mixed speeds), $\sigma \le 0.35$ & Urgency-weighted \cb & \texttt{share\_map=true} on every \texttt{BID}; \texttt{refuge\_blocked} enabled & $0.04$\,s weighted excess ($98\%$ optimal match vs.\ $27.1$\,s fixed priority); ahead of both rulebooks up to $\sigma=0.35$, tied at $0.4$--$0.5$; avoids Lidar-lag penalty of \texttt{share\_on\_detect} \\\addlinespace[4.5pt]
Self-interested / multi-vendor robots sharing aisles & \texttt{clarke} (VCG pivot payment to operator ledger) & Authenticated IDs \& bay sets; audit joint bid scaling & $0.00$\,s unilateral misreport gain ($\le 0.02$\,s realized), with shared-snapshot types; generates a $24$--$37\%$ bank surplus (of realized delay) and $+4.9$\,s collusion incentive \\\addlinespace[4.5pt]
Convoys (2-vs.-1) \& branch junctions & Convoys: \texttt{fixed\_priority}; junctions: \cb; all with \texttt{make\_room}/\texttt{goal\_wait} & Fuse fresh \texttt{STATE} with Lidar for occluded junction lanes & \texttt{make\_room} \& \texttt{goal\_wait} restore $100\%$ completion; urgency rulebook is worst in convoys ($133.6$ vs.\ $90.3$\,s at $u_B=3$); model \emph{serial} follower clearance when $u_B \gg 1$ \\\addlinespace[4.5pt]
Legacy aisle without pull-in bays or off-lane staging queues & \texttt{reservation} (only if staging sits off exit path) & Corridor entry/exit \texttt{STATE} beacons & $100\%$ deadlock-free, but costs $+18.8$\,s excess when mid-aisle bays exist and staging blocks peer exit \\\bottomrule
\end{tabular}}
\end{table*}

Three broader lessons emerge from the experiments. First, \emph{broadcasting snapshots with attached bay lists outperforms complex handshakes}. Periodically broadcasting \texttt{BID} snapshots that include known and blocked bay IDs eliminates incomplete-map penalties at zero extra delay. By contrast, waiting for an explicit \texttt{commit} confirmation adds $3$--$4$\,s of delay under extreme packet loss without reducing collisions or deadlocks, and waiting for Lidar peer detection before sharing maps (\texttt{share\_on\_detect}) delivered no peer map before the decision in 3D Gazebo (0/6 runs) because of point-cloud clustering lag, sending one of the six runs to the far end room. Second, \emph{the benefit of cost bidding depends on matching the right fleet objective and having low estimation noise}. Minimizing unweighted delay when robots have different urgency weights wastes twice as much weighted time ($7.0$\,s) as a simple six-rule heuristic ($3.5$\,s), and, in the single-bay layout, bid estimation noise of $\sigma=0.5$ erases the advantage of \cb\ over a closest-refuge rulebook (which itself does not carry over to layouts with several bays). Third, \emph{simplified multi-robot delay models break down before 2D kinematics do}. In two-robot straight corridors, 2D excess delays transfer to 3D Gazebo within $1.0$\,s. However, three-robot convoys and T-junctions introduce sequential backing waits and branch-exit timing races that simple additive delay sums miss, requiring multi-robot timeline modeling and 3D physics validation.
\subsection{Limitations and Future Scope}
This study has several clear boundaries that point to directions for future work.
\begin{itemize}
\item \emph{Simulation-only evaluation and ground-truth localization.} All experiments were conducted in 2D and Gazebo Harmonic simulation (using 6 to 10 runs per Gazebo condition). Because the Gazebo nodes read ground-truth robot pose to isolate negotiation and Lidar perception, real-world factors such as AMCL or SLAM localization drift, wheel slip, and physical TurtleBot3 hardware dynamics remain to be tested.
\item \emph{Matched motion model and non-bid cheating.} Both \cb\ and \texttt{clarke} calculate bids using the simulator's true speed limits and bay coordinates. Systematic model errors (such as miscalibrated reverse speeds), strategic bay hiding, robot-ID spoofing during \texttt{commit} tie-breaks, mechanisms for redistributing accumulated VCG ledger credits, and voluntary participation incentives are not evaluated here.
\item \emph{Reactive baseline and theoretical liveness guarantees.} The only zero-communication baseline tested was centerline \texttt{reactive} backoff rather than continuous local planners such as ORCA, RVO, or CBFs. In addition, the protocol does not yet have a formal mathematical bound on how quickly temporary mutual yields ($14\%$--$24\%$ of runs under extreme packet loss) resolve, and the best forced yielder serves as a decision-level reference rather than a globally optimal continuous trajectory.
\item \emph{Simplified layouts and decision latency.} The T-junction approximates a side corridor as a $2$\,m-deep bay (leaving an unexplained $-345.4$\,s cumulative delay gap between Gazebo and 2D on the turn-in variant), \texttt{explore} discovers bays by geometric radius rather than full SLAM, the map-mismatch and blocked-bay sweeps use 10 seeds per condition, and the \texttt{learned} policy is a centralized tabular outcome classifier rather than deep MARL. Finally, every policy evaluated in this paper computes its decision in less than a millisecond of CPU time. Using large language or vision-language models (LLMs/VLMs)~\cite{marlin2024,gamechat2025,garg2024llm} as negotiators introduces multi-second inference delays while robots are still moving, along with decentralized disagreement failures that require dedicated asynchronous protocols and separate investigation.
\end{itemize}

